% vim: spelllang=en ts=2 sw=2 sts=2
\documentclass[11pt,a4paper]{scrartcl}

\usepackage[utf8]{inputenc}
\usepackage[T1]{fontenc}
\usepackage{amsmath}
\usepackage{amssymb}
\usepackage{amsthm}
\usepackage{mathtools}
\usepackage{thmtools}
\usepackage{enumitem}
\setlist[enumerate]{label=(\alph*),left=0pt,labelsep=1em}
\usepackage{multicol}
\usepackage{array}
\usepackage{xcolor}



%%%%%%%%%%%%%%%%%%%%%%%Format%%%%%%%%%%%%%%%%%%%%%%%
\setlist[description]{labelindent=\parindent}
\usepackage[left=2cm,right=2cm,top=2cm,bottom=2.5cm]{geometry}

\KOMAoptions{parskip=never}


\AtBeginDocument{
	\setlength{\abovedisplayskip}{5pt}
	\setlength{\belowdisplayskip}{5pt}
	\setlength{\abovedisplayshortskip}{2pt}
	\setlength{\belowdisplayshortskip}{2pt}
	\setlength{\jot}{2pt}
}


\usepackage{hyperref}
\hypersetup{
	colorlinks,
	linkcolor={magenta!70!black},
	citecolor={green!60!black},
	urlcolor={blue!60!black}
}
\usepackage[capitalise]{cleveref}

\addtokomafont{disposition}{\rmfamily}  %same font everywhere

\setcounter{tocdepth}{2}

%Fix spacing around "\left(\right)"
\usepackage{mleftright} 
\mleftright % redefine \left as \mleft and \right as \mright.
%%%%%%%%%%%%%%%%%%%%%%%%%%%%%%%%%%%%%%%%%%%%%%%%%%%%%

%Bibliography
\usepackage[
backend=biber,
style=alphabetic,
sorting=nyt,
isbn=false
]{biblatex}
\renewbibmacro*{doi+eprint+url}{%if DOI is available, no URL
	\iftoggle{bbx:doi}
	{\printfield{doi}}
	{}%
	\newunit\newblock
	\ifboolexpr{togl {bbx:eprint} and test {\iffieldundef{doi}}}
	{\usebibmacro{eprint}}
	{}%
	\newunit\newblock
	\ifboolexpr{togl {bbx:url} and test {\iffieldundef{doi}}  and test {\iffieldundef{eprint}}}
	{\usebibmacro{url+urldate}}
	{}}
\addbibresource{bibliography.bib}


%Diagrams
\usepackage{tikz}
\usetikzlibrary{calc}
%\usepackage{tikzit/tikzit}
%\input{tikzit/markov2.tikzstyles}
\usepackage{tikz-cd}


%draft packages
%\usepackage{showkeys}
%\usepackage{todonotes}
%\newcommand{\note}[1]{\todo[color=blue!20,inline,caption={}]{#1}}
%\newcommand\myworries[1]{\textcolor{red}{\textbf{Note:} #1}}


%%%%%%%%%%%%%%%%%%%%%%%Commands%%%%%%%%%%%%%%%%%%%%%%%

%categories
\newcommand{\cat}{\mathsf{C}}
\newcommand{\catd}{\mathsf{D}}
\newcommand{\setc}{\mathsf{Set}}
\newcommand{\finset}{\mathsf{FinSet}}
\newcommand{\cntblset}{\mathsf{CntblSet}}
\newcommand{\setmulti}{\mathsf{SetMulti}}
\newcommand{\chaus}{\mathsf{CHaus}}
\newcommand{\meas}{\mathsf{Meas}}
\newcommand{\borelmeas}{\mathsf{BorelMeas}}
\newcommand{\topc}{\mathsf{Top}}
\newcommand{\nat}{\mathsf{Nat}}
\newcommand{\finstoch}{\mathsf{FinStoch}}

\newcommand{\simp}{\mathsf{s}}

%numbers
\newcommand{\N}{\mathbb{N}}
\newcommand{\Z}{\mathbb{Z}}
\newcommand{\Q}{\mathbb{Q}}
\newcommand{\R}{\mathbb{R}}
\newcommand{\C}{\mathbb{C}}

%misc 
\newcommand{\id}{\operatorname{id}}
\newcommand{\im}{\operatorname{im}}
\newcommand{\dx}{\mathrm{d}}
\newcommand{\Hom}{\operatorname{Hom}}
\newcommand{\Obj}{\operatorname{Obj}}
\newcommand{\cop}{\operatorname{copy}}
\newcommand{\del}{\operatorname{del}}
\newcommand{\detc}[1]{#1_{\operatorname{det}}}
\newcommand{\dec}{\operatorname{Dec}}

\newcommand{\power}{\mathcal{P}}
\newcommand{\finiteset}[1]{[#1]}	% finite set {1, .., n} or perhaps {0, .., n-1}, depending on which convention is more convenient
\newcommand{\one}{\finiteset{1}}		% one-element set
\newcommand{\two}{\finiteset{2}}		% two-element set
\newcommand{\three}{\finiteset{3}}		% three-element set

\newcommand{\doubleslash}{/\!\!/}
\newcommand{\dist}{\mathcal{D}}
\newcommand{\distfin}{\mathcal{D}_{\mathrm{fin}}}
\newcommand{\giry}{\mathcal{P}}
\newcommand{\Sym}{\operatorname{Sym}}




%%%%%%%%%%%%%%%%%%%%%%%%%%%%%%%%%%%%%%%%%%%%%%%%%%%%%%



%%%%%%%%%%%%%%%%%%%%%%%Theorems%%%%%%%%%%%%%%%%%%%%%%%

\newtheorem{theo}{Theorem}[section]
\newtheorem{defi}[theo]{Definition}
\newtheorem{prop}[theo]{Proposition}
\newtheorem{lem}[theo]{Lemma}
\newtheorem{cor}[theo]{Corollary}
\newtheorem{conj}[theo]{Conjecture}
\newtheorem{prob}[theo]{Problem}
\newtheorem{que}[theo]{Question}

\theoremstyle{definition}
\newtheorem{rem}[theo]{Remark}
\newtheorem*{rem*}{Remark}
\newtheorem{ex}[theo]{Example}




%cleveref
\Crefname{theo}{Theorem}{Theorems}
\Crefname{defi}{Definition}{Definitions}
\Crefname{prop}{Proposition}{Propositions}
\Crefname{lem}{Lemma}{Lemmas}
\Crefname{cor}{Corollary}{Corollaries}
\Crefname{conj}{Conjecture}{Conjectures}
\Crefname{rem}{Remark}{Remarks}
\Crefname{ex}{Example}{Examples}
\Crefname{equation}{}{}		% write '(2.3)' instead of 'Equation (2.3)'
\Crefname{figure}{Figure}{Figures}







\title{Naturality characterizes product measures and relative product measures}
%\subtitle{}

\author{Victor Bloch\thanks{Supported by a doctoral scholarship of the University of Innsbruck.}, Tobias Fritz\\[0.5em]
\small Department of Mathematics, University of Innsbruck}

\date{\today}


\begin{document}

\maketitle

\begin{abstract}
	We show that the product measure is the only \emph{natural} way to assign to each pair of probability measures on measurable spaces a probability measure on their product.
	Here, naturality is meant in the sense of category theory and amounts to the condition that the assignment commutes with pushforward along measurable maps.
	In the standard Borel setting, we also prove an analogous result for relative product measures over a fixed base probability space.
\end{abstract}

\tableofcontents

\section{Introduction}

Categorical approaches to probability theory and statistics trace back to Lawvere's seminal 1962 manuscript on the category of probabilistic mappings~\cite{Lawvere1962} and Giry's 1982 paper on the monad of probability measures~\cite{Giry1982}.
But especially in recent years, there has been a surge of interest with surprisingly strong results.
To name a few, these include fully categorical proofs of de Finetti's theorem~\cite{Fritz2018}, the Aldous--Hoover theorem~\cite{Chen2025}, the $d$-separation criterion~\cite{FritzKlingler2023}
and categorical accounts of ergodic decomposition~\cite{Moss2023}, the law of large numbers~\cite{Fritz2026} and the Metropolis--Hastings algorithm~\cite{Cornish2026}.
A recurring theme is that probabilistic structures are abstractly encoded in categorical terms, and this leads to a higher-level understanding of their properties and interrelations,
and sometimes also to intriguing new questions.

In the present work, we raise and answer a question of this kind: what singles out the product-measure construction---and its generalization to relative products over a fixed base probability space---among all possible ways of combining two probability measures into a joint distribution?
The simple answer is that the categorical notion of \emph{naturality} completely characterizes it.
In measure-theoretic terms, naturality means that the measure associated to a pair of pushforward measures is the pushforward of the measure associated to the original pair.
We show in \cref{theo:onenaturaltransformation} that taking product measures is the only way to achieve this.
This is arguably not very surprising, and indeed also closely related to a result in copula theory that we discuss below.

Product measures are a fundamental concept in probability theory associated with stochastic independence.
The more involved concept of \emph{conditional independence} is likewise related to the construction of \emph{relative product measures} over a fixed base probability space,
where the measure on this base space may be thought of as formalizing shared information between the two spaces above it.
As is the case in categorical approaches to topology and geometry as well, such relative constructions are particularly amenable to a categorical treatment via slice categories.
In fact, we will extend \cref{theo:onenaturaltransformation} to a relative version in \cref{theo:onenaturaltransformation_relative}.
It characterizes the relative product measure as the uniquely natural way of combining two probability measures over a fixed base probability space into a joint distribution.
As the formation of relative product measures itself, this characterization requires stronger assumptions on the measurable spaces involved, and we only consider the case of standard Borel spaces.

\subsection*{Related work}

The present work is closely related to \emph{copula theory},
which studies systematic ways of combining marginal distributions into joint distributions~\cite{Nelsen2006,Durante2007,Durante2009}.
If one focuses on atomless distributions on standard Borel spaces,
then one can restrict to the Lebesgue measure on $[0,1]$ as marginals without loss of generality.
Hence an $n$-variate copula is defined as a probability measure on $[0,1]^n$ with uniform marginals.
For example, the \emph{independence copula} is the product measure on $[0,1]^n$.

Close to our results is the characterization of the independence copula as the unique copula that is invariant under all measure-preserving transformations of its coordinates due to Durante, Sarkoci and Sempi~\cite[Theorem~3.1]{Durante2009}.
In fact, this can be viewed as a version of our \cref{theo:onenaturaltransformation},
namely the restriction to uncountable standard Borel spaces and considering the functor which maps every such space to the non-atomic probability measures on it.
The proof amounts to noting that the Lebesgue measure on $[0,1]^2$ is the only measure that assigns the same number to every smaller square in an $n \times n$ partition.

Some of the arguments in this manuscript can be formulated and proved elegantly in terms of the abstract framework of Markov categories~\cite{Fritz2020}.
However, we prefer to spell out the measure-theoretic details here in order to avoid the additional overhead of recalling the Markov-categorical framework,
given that the main results are about the classical functors of probability measures only.

\section{Preliminaries}

\subsection{Basics}

In this subsection, we declare some notation and recall some standard material on functors of probability measures.
For natural transformations, we usually omit the subscripts which index the components, as these can be inferred from context.
For $n \in \N$, we write $\finiteset{n} \coloneqq \{1,\ldots,n\}$.

Denote by $\setc$ the category of sets and functions. Let $\dist: \setc \rightarrow \setc$ be the \emph{distribution functor}. It sends a set $X$ to the set $\dist X$ of discrete probability measures on $X$. A discrete probability measure is a finitely supported function $p: X \rightarrow [0,1]$ s.t. $\sum_{x \in X} p(x) = 1$, 
or equivalently a probability measure on the power set $\power(X)$ as $\sigma$-algebra.
On functions $f : X \to Y$, $\dist$ acts by taking the pushforward, $(Df)(p) \coloneq f_\ast p$, where $(f_\ast p)(S) = p(f^{-1}(S))$ for all $S \subseteq Y$.
Functoriality of $\dist$ is straightforward to check.
We also write $\distfin : \finset \rightarrow \setc$ for the restriction of $\dist$ to $\finset$, the full subcategory of $\setc$ consisting of finite sets.

The most general framework for measure-theoretic probability is provided by $\meas$, the category of measurable spaces and measurable maps. A probability measure on a measurable space $(X,\Sigma_X)$ is a function $p: \Sigma_X \rightarrow [0,1]$ that satisfies $\sigma$-additivity and the normalization equation $p(X) = 1$.
For a measurable map $f:(X,\Sigma_X) \rightarrow (Y,\Sigma_Y)$, the pushforward $f_\ast p$ is again defined by $(f_\ast p)(S) \coloneq p(f^{-1}(S))$ for $S \in \Sigma_Y$. The analog of $\dist$ now is the \emph{Giry functor} $\giry: \meas \rightarrow \setc$~\cite{Giry1982}.
This functor sends a measurable space $(X,\Sigma_X)$ to the set of all probability measures on it.\footnote{As for the Giry monad~\cite{Giry1982}, we could also endow every $\giry X$ with the $\sigma$-algebra generated by the evaluation maps and therefore make $\giry$ land in $\meas$ itself,
but we will not need this here.}
On measurable maps, $\giry$ acts as the pushforward, that is, for $f: (X,\Sigma_X) \rightarrow (Y,\Sigma_Y)$ and $p$ a probability measure on $(X,\Sigma_X)$ we have $(\giry f (p) \coloneq f_\ast p$.

A measurable space $(X,\Sigma_X)$ is a \emph{standard Borel space} if there is a separable and completely metrizable topology on $X$ such that $\Sigma_X$ is the $\sigma$-algebra generated by its open sets.
Standard Borel spaces play a fundamental role in probability theory, as they are both very well-behaved and sufficiently general to include most measurable spaces of interest.
We denote the full subcategory of $\meas$ on standard Borel spaces by $\borelmeas$.
The Giry functor restricts to an endofunctor of $\borelmeas$, and we also write $\giry: \borelmeas \rightarrow \setc$ for this restriction.

\subsection{Product measures}

Given two measurable spaces $(X,\Sigma_X)$ and $(Y,\Sigma_Y)$, we endow $X \times Y$ with the \emph{product $\sigma$-algebra} $\Sigma_X \otimes \Sigma_Y$,
which is generated by the rectangles $S \times T$ with $S \in \Sigma_X$ and $T \in \Sigma_Y$.
This is the categorical product both on $\meas$ and on $\borelmeas$ with respect to the usual projection maps
\[
	  \pi_X \: : \: X \times Y \longrightarrow X, \qquad
	  \pi_Y \: : \: X \times Y \longrightarrow Y.
\]
It clearly restricts to the usual categorical product on $\setc$ and $\finset$.

For probability measures~$p$ on $(X,\Sigma_X)$ and $q$ on $(Y,\Sigma_Y)$, a \emph{joint distribution} is a probability measure $r$ on $(X \times Y, \Sigma_X \otimes \Sigma_Y)$ such that $p$ and $q$ are the marginals of $r$, that is, 
\[
	(\pi_X)_\ast r = p, \qquad (\pi_Y)_\ast r = q.
\]
There are typically many such joint distributions $r$, which differ in how they correlate the two marginals $p$ and $q$.
The canonical choice, corresponding to no correlation, is the product measure:

\begin{prop}[Product measure] \label{defi:productmeasure}
  Let $(X,\Sigma_X)$ and $(Y,\Sigma_Y)$ be measurable spaces and $p \in \giry(X)$ and $q \in \giry(Y)$.
  Then there is a unique measure $p \otimes q$ on $(X \times Y, \Sigma_X \otimes \Sigma_Y)$ such that
  \begin{equation}
    \label{eq:productmeasure}
    (p \otimes q) (S \times T) = p(S) \cdot q(T)
  \end{equation}
  for all $S \in \Sigma_X$ and $T \in \Sigma_Y$.
\end{prop}

The formation of product measures defines maps
\begin{align*}
	\nabla \: : \: \giry X \times \giry Y & \longrightarrow \giry(X \times Y) \\
	(p,q) & \longmapsto p \otimes q
\end{align*}
which are natural in $X$ and $Y$, meaning that for all measurable $f : X \to X'$ and $g : Y \to Y'$, we have
\[
	(f \times g)_\ast (p \otimes q) = f_\ast p \otimes g_\ast q.
\]
To see this, we evaluate on any rectangle $S' \times T'$ with $S' \in \Sigma_{X'}$ and $T' \in \Sigma_{Y'}$, where we get
\begin{align*}
	(f_\ast p \otimes g_* q)(S' \times T') & = (f_\ast p)(S') \cdot (g_\ast q)(T') \\[2pt]
						& = p(f^{-1}(S')) \cdot q(g^{-1}(T')) \\[2pt]
						& = (p \otimes q)(f^{-1}(S') \times g^{-1}(T')) \\[2pt]
						& = ((f \times g)_\ast (p \otimes q))(S' \times T'),
\end{align*}
and this implies the desired equality by the uniqueness part of \cref{defi:productmeasure}.
Hence the formation of product measures can be seen as a natural transformation
\begin{equation}
	\label{eq:productnaturality_giry}
	\nabla \: : \: \giry\_ \times \giry\_ \Longrightarrow \giry(\_ \times \_),
\end{equation}
where both sides are functors $\meas \times \meas \longrightarrow \setc$.\footnote{In fact, this natural transformation makes $\giry$ into a lax symmetric monoidal functor~\cite{Fritz2018}.}
All of this applies likewise to the distribution functor $\dist$ on $\setc$: the formation of product distributions is a natural transformation
\begin{equation}
	\label{eq:productnaturality_dist}
	\nabla \: : \: \dist\_ \times \dist\_ \Longrightarrow \dist(\_ \times \_)
\end{equation}
between functors $\setc \times \setc \longrightarrow \setc$, and likewise for $\distfin$ on $\finset$.

\subsection{Relative product measures}

Product measures are about probabilistic independence: if one views $X$ and $Y$ as random variables on the probability space $(X \times Y, p \otimes q)$ via the projection maps,
then $X$ and $Y$ are independent, and the product measure is the unique joint distribution with this property.
But even more important is the concept of \emph{conditional independence},
and the relevant notion of product measure is the \emph{relative product measure}~\cite{Dawid1999},
which we now review.
These relative products have been studied under various names in different areas,
e.g.~in the context of graphical models~\cite{Dawid1999} and ergodic theory~\cite[Definition~6.15]{Einsiedler2011}.
Categorical treatments include~\cite{Flori2016,Simpson2018,Perrone2025} and~\cite[Definition~12.8]{Fritz2020}.\footnote{These works partly restrict to the case where $f$ and $g$ are themselves product projections, but modulo adding outcomes of measure zero, this is no restriction.}

Similarly to how product measures live on product spaces, relative product measures live on pullbacks.
So let
\begin{equation}
	\label{eq:pullbackdiagram}
	\begin{tikzcd}
		X \times_B Y \arrow[r, "\pi_X"] \arrow[d, "\pi_Y"'] & X \arrow[d, "\rho_X"] \\
		Y \arrow[r, "\rho_Y"'] & B
	\end{tikzcd}
\end{equation}
be a pullback diagram in $\meas$.
Suppose that $p \in \giry(X)$ and $q \in \giry(Y)$ are probability measures such that $(\rho_X)_\ast p = (\rho_Y)_\ast q$.
Then the question is: is there a canonical probability measure $r \in \giry(X \times_B Y)$ such that $(\pi_X)_\ast r = p$ and $(\pi_Y)_\ast r = q$?

If $X$, $Y$ and $B$ are standard Borel,
and if we view $X$, $Y$ and $B$ as random variables on the probability space $(X \times_B Y, r)$ via the given maps,
then there is a unique $r$ which makes $X$ and $Y$ conditionally independent given $B$~\cite[Proposition~12.9]{Fritz2020}.
This is the relative product of $p$ and $q$ over $B$, and we denote it by
\[
	p \otimes_\omega q \:\in\: \giry(X \times_B Y),
\]
where $\omega \coloneq (\rho_X)_\ast p = (\rho_Y)_\ast q$ is the common pushforward on $B$.
Writing down $p \otimes_\omega q$ explicitly requires regular conditional probabilities.
These are most easily understood as Markov kernels, which we now recall.

\begin{defi}
	For measurable spaces $(X, \Sigma_X)$ and $(Y, \Sigma_Y)$, a \emph{Markov kernel} from $X$ to $Y$ is a map
	\begin{align*}
		\kappa \: : \: \Sigma_Y \times X & \longrightarrow [0,1] \\
		(S, x) & \longmapsto \kappa(S|x)
	\end{align*}
	such that:
	\begin{enumerate}[label=(\roman*)]
		\item $\kappa(-|x)$ is a probability measure on $(Y, \Sigma_Y)$ for all $x \in X$.
		\item $\kappa(S|\cdot)$ is measurable for all $S \in \Sigma_Y$.
	\end{enumerate}
\end{defi}

Equivalently, a Markov kernel can be viewed as a measurable map $X \to \giry Y$.\footnote{The equivalence of these two points of view leads to the definition of \emph{representable Markov category}~\cite[Definition~3.10]{Fritz2023}.}

\begin{defi}
	Let $\rho_X : X \to B$ be a measurable map between measurable spaces $(X,\Sigma_X)$ and $(B,\Sigma_B)$.
	Then a \emph{regular conditional probability} of a probability measure $p \in \giry X$ with respect to $\rho_X$ is a Markov kernel $p(-|\cdot) : \Sigma_X \times B \to [0,1]$ such that
  \begin{equation}
		\label{eq:disintegration}
		p(S \cap \rho_X^{-1}(T)) = \int_{b \in T} p(S|b) \  \omega(\dx b)
		\qquad
		\forall S \in \Sigma_X, T \in \Sigma_B,
	\end{equation}
	where $\omega \coloneq (\rho_X)_\ast p$.
\end{defi}

\begin{prop}\label{prop:regularconditionalprobability}
	If $X$ and $B$ are standard Borel, then such a regular conditional probability exists for every $p \in \giry X$,
	and it is unique up to $\omega$-a.e.~equality.
\end{prop}

See e.g.~\cite[Lemma 10.4.3 \& Proposition 10.4.12]{Bogachev2007} for proofs of this standard result.
To understand the definition, it may also help to note that
\begin{equation}
	\label{eq:fiberconcentration}
	p(\rho_X^{-1}(b)|b) = 1
\end{equation}
for $\omega$-almost all $b \in B$~\cite[Example~10.4.11]{Bogachev2007}.
In other words, $p(-|\cdot)$ is essentially a stochastic section of $\rho_X$.

\begin{defi}
  Let $X$, $Y$, and $B$ be standard Borel spaces with measurable maps $\rho_X : X \rightarrow B$ and $\rho_Y : Y \rightarrow B$,
	and fix $\omega \in \giry B$.
	Suppose that $p \in \giry X$ and $q \in \giry Y$ satisfy $(\rho_X)_\ast p = \omega = (\rho_Y)_\ast q$.
	Then the \emph{relative product} of $p$ and $q$ over $\omega$ is the probability measure on $X \times Y$ defined by
  \begin{equation}
		\label{eq:relativeproductmeasure}
		(p \otimes_\omega q)(S) \coloneqq \int_{b \in B} \left(p(-|b) \otimes q(-|b)\right)(S) \ \omega(\dx b)
		\qquad
		\forall S \in \Sigma_{X \times Y}.
  \end{equation}
\end{defi}

Due to the uniqueness of regular conditional probabilities up to $\omega$-a.e.~equality, the integral and therefore the relative product measure are well-defined;
the monotone convergence theorem implies $\sigma$-additivity.

If $B$ is finite, then the relative product measure is characterized by
\[
	(p \otimes_\omega q)(S \times T) = \sum_{b \in B \: : \: \omega(b) > 0} \frac{p(S \cap \rho_X^{-1}(b)) \cdot q(T \cap \rho_Y^{-1}(b))}{\omega(b)}.
\]
In particular if $X$ and $Y$ are also discrete, then the probabilities of individual outcomes are given by
\[
	(p \otimes_\omega q)(x,y) = \begin{cases}
		\frac{p(x) \cdot q(y)}{\omega(\rho_X(x))} = \frac{p(x) \cdot p(y)}{\omega(\rho_Y(y))} & \text{if } \rho_X(x) = \rho_Y(y) \text{ and } \omega(\rho_X(x)) > 0, \\
		0 & \text{otherwise.}
	\end{cases}
\]
The zero case is a manifestation of the following general support property.

\begin{lem}[{e.g.~\cite[Prop.~6.16(1)]{Einsiedler2011}}]
	\label{lem:relativeproductsupport}
	The relative product measure $p \otimes_\omega q$ is supported on the pullback
	\begin{equation}
		\label{eq:pullbackdiagram2}
		\begin{tikzcd}
			X \times_B Y \arrow[r, "\pi_X"] \arrow[d, "\pi_Y"'] & X \arrow[d, "\rho_X"] \\
			Y \arrow[r, "\rho_Y"'] & B
		\end{tikzcd}
	\end{equation}
	considered as a measurable subset of $X \times Y$.
\end{lem}

\begin{proof}
	Since $B$ is standard Borel, its diagonal $\Delta_B \subseteq B \times B$ is measurable. Hence
	\[
		X \times_B Y = (\rho_X \times \rho_Y)^{-1}(\Delta_B)
	\]
	is indeed a measurable subset of $X \times Y$.
	Moreover, the fiber concentration property~\eqref{eq:fiberconcentration} gives
	\[
		p(\rho_X^{-1}(b)|b) = 1,
		\qquad
		q(\rho_Y^{-1}(b)|b) = 1
	\]
	for $\omega$-almost every $b \in B$.
	For each such $b$, we have
	\[
		\rho_X^{-1}(b) \times \rho_Y^{-1}(b) \subseteq X \times_B Y,
	\]
	and therefore
	\[
		\bigl(p(-|b) \otimes q(-|b)\bigr)(X \times_B Y) = 1.
	\]
	Thus with $S = X \times_B Y$ in~\eqref{eq:relativeproductmeasure}, the integrand in~\eqref{eq:relativeproductmeasure} is a.e.~equal to $1$,
	and hence the integral is $1$ as well.
\end{proof}

We now turn to further categorical aspects of relative product measures.
Given that the construction takes place ``over $B$'', it is natural to work with slice categories like $\borelmeas/B$.
An object of $\borelmeas/B$ is a pair $(X, \rho_X)$ where $X$ is a standard Borel space and $\rho_X: X \rightarrow B$ is a measurable map.
We often suppress the map $\rho_X$ in the notation and simply write $X$, thinking of $\rho_X$ as extra structure on the measurable space $X$.
In this slice category, the pullback diagram~\eqref{eq:pullbackdiagram2} is the categorical product of $X$ and $Y$ over $B$.

We also fix the base probability measure $\omega \in \giry B$ throughout.
Then the relative version of the probability measures functor is
\begin{align*}
	\giry^\omega \: : \: \borelmeas/B & \longrightarrow \borelmeas \\
	(X, \rho_X) & \longmapsto \{ p \in \giry X \: : \: (\rho_X)_\ast p = \omega \}.
\end{align*}
The action on a morphism of $\borelmeas/B$, which is a commutative triangle of the form
\[
	\begin{tikzcd}
		X \arrow[rr, "f"] \arrow[dr, "\rho_X"'] & & Y \arrow[dl, "\rho_Y"] \\
		& B &
	\end{tikzcd}
\]
is again given by the pushforward $f_\ast$.

Our next goal is to show that the formation of relative product measures defines a natural transformation,
but this requires a bit of preparation.

\begin{lem}\label{lem:pushforwarddisintegration}
  Let $f : X \to Y$ be a morphism in $\borelmeas/B$.
	Then for every $p \in \giry^\omega X$, the Markov kernel
	\begin{align*}
		B & \longrightarrow \giry Y \\
		b & \longmapsto f_\ast p(-|b)
	\end{align*}
	is a regular conditional probability of $f_\ast p$ with respect to $\rho_Y$.
\end{lem}

\begin{proof}
	The displayed map is a Markov kernel because it is the composite of the Markov kernel
	$p(-|\cdot) : B \to \giry X$ with the measurable pushforward map $f_\ast : \giry X \to \giry Y$.
	Since $f$ is a morphism over $B$, we have
	\[
		(\rho_Y)_\ast(f_\ast p)
		= (\rho_Y \circ f)_\ast p
		= (\rho_X)_\ast p
		= \omega.
	\]
	Consequently, for $S \in \Sigma_Y$ and $T \in \Sigma_B$,
	\begin{align*}
		(f_\ast p)(S \cap \rho_Y^{-1}(T))
		&= p\bigl(f^{-1}(S) \cap \rho_X^{-1}(T)\bigr) \\
		&= \int_{b \in T} p(f^{-1}(S)|b) \, \omega(\dx b) \\
		&= \int_{b \in T} \bigl(f_\ast p(-|b)\bigr)(S) \, \omega(\dx b),
	\end{align*}
	which is precisely the disintegration identity~\eqref{eq:disintegration}.
\end{proof}

Here is now the promised naturality statement for relative product measures.

\begin{lem}
	\label{lem:relativeproductnaturality}
	The formation of relative product measures defines maps
	\[
		\nabla^\omega \: : \: \giry^\omega X \times \giry^\omega Y \longrightarrow \giry^\omega(X \times_B Y)
	\]
	that are natural in $X, Y \in \borelmeas/B$.
\end{lem}

For the distribution functor $\distfin$ in place of $\giry$, this statement is also part of~\cite[Proposition~5.2.1]{Flori2016}.

\begin{proof}
	Let $p \in \giry^\omega X$ and $q \in \giry^\omega Y$.
	By \cref{lem:relativeproductsupport}, we can regard $p \otimes_\omega q$ as a probability measure on $X \times_B Y$.
	Its first marginal is $p$, since for every $S \in \Sigma_X$,
	\begin{align*}
		((\pi_X)_\ast(p \otimes_\omega q))(S)
		&= \int_{b \in B} p(S|b)q(Y|b) \, \omega(\dx b) \\
		&= \int_{b \in B} p(S|b) \, \omega(\dx b) \\
		&= p(S).
	\end{align*}
	It follows that
	\[
		(\rho_X \circ \pi_X)_\ast(p \otimes_\omega q)
		= (\rho_X)_\ast p
		= \omega,
	\]
	so $p \otimes_\omega q \in \giry^\omega(X \times_B Y)$, making $\nabla^\omega$ well-defined.

	Suppose that we are given two morphisms in $\borelmeas/B$,
	\[
		\begin{tikzcd}[column sep=large]
			X \arrow[rr, "f"] \arrow[dr, "\rho_X"'] & & X' \arrow[dl, "\rho_{X'}"] \\
			& B &
		\end{tikzcd}
		\qquad\quad
		\begin{tikzcd}[column sep=large]
			Y \arrow[rr, "g"] \arrow[dr, "\rho_Y"'] & & Y' \arrow[dl, "\rho_{Y'}"] \\
			& B &
		\end{tikzcd}
	\]
	By \cref{lem:pushforwarddisintegration}, regular conditional probabilities of $f_\ast p$ and $g_\ast q$ are given by
	$b \mapsto f_\ast p(-|b)$ and $b \mapsto g_\ast q(-|b)$, respectively.
	Therefore, for every $S \in \Sigma_{X' \times Y'}$, naturality of ordinary product measures gives
	\begin{align*}
		(f_\ast p \otimes_\omega g_\ast q)(S)
		&= \int_{b \in B} \bigl(f_\ast p(-|b) \otimes g_\ast q(-|b)\bigr)(S) \, \omega(\dx b) \\
		&= \int_{b \in B} \bigl(p(-|b) \otimes q(-|b)\bigr)((f \times g)^{-1}(S)) \, \omega(\dx b) \\
		&= (p \otimes_\omega q)((f \times g)^{-1}(S)) \\
		&= \bigl((f \times g)_\ast(p \otimes_\omega q)\bigr)(S).
	\end{align*}
	Since $f$ and $g$ are morphisms over $B$, the map $f \times g$ restricts to
	$f \times_B g : X \times_B Y \to X' \times_B Y'$.
	Together with \cref{lem:relativeproductsupport}, the displayed equality gives
	\[
		(f \times_B g)_\ast(p \otimes_\omega q)
		= f_\ast p \otimes_\omega g_\ast q,
	\]
	which is the desired naturality.
\end{proof}

So similar to the formation of product measures,
the formation of relative product measures can be seen as a natural transformation
\begin{equation}
	\label{eq:relativeproductnaturality_giry}
	\nabla^\omega \: : \: \giry^\omega \_ \times \giry^\omega \_ \Longrightarrow \giry^\omega(\_ \times_B \_),
\end{equation}
where both sides are functors $\borelmeas/B \times \borelmeas/B \longrightarrow \setc$.

\begin{rem}
	\begin{enumerate}
		\item 
			If $B$ is a singleton, then the relative product reduces to the ordinary product,
			and \cref{lem:relativeproductnaturality} reduces to the naturality of ordinary product measures.
		\item 
			Clearly all of this applies likewise to the distribution functor $\dist$ on $\setc$:
			for any set $B$ and $\omega \in \dist B$,
			the formation of relative product distributions is a natural transformation
			\begin{equation}
				\label{eq:relativeproductnaturality_dist}
				\nabla^\omega \: : \: \dist\_ \times \dist\_ \Longrightarrow \dist(\_ \times \_)
			\end{equation}
			between functors $\setc/B \times \setc/B \longrightarrow \setc$.
	\end{enumerate}
\end{rem}

%\begin{rem}
%	In fact, the relative product also equips $\giry^\omega$ with the structure of a lax symmetric monoidal functor.
%	\note{Das kann man sogar aus der Eindeutigkeit folgern, falls wir sie nicht nur für zwei sondern für $n$ Faktoren zeigen.}
%	We leave the straightforward details to the reader.
%	For finite sets, the relevant associativity statement can also be seen as part of~\cite[Proposition~5.2.1]{Flori2016}.
%\end{rem}

\section{Main results}

\subsection{Statements}

A natural question now is, are there any other natural transformations of the form~\eqref{eq:productnaturality_giry} or \eqref{eq:productnaturality_dist}? 
Or is the formation of product measures the only natural way to combine two probability measures into a joint measure?
This is the first question that we answer in the affirmative.

\begin{theo}
	\label{theo:onenaturaltransformation}
	\begin{enumerate}
		\item 
			With $\giry : \meas \to \setc$ the Giry functor, the formation of product measures is the only family of maps
			\[
				\giry X \times \giry Y \longrightarrow \giry(X \times Y)
			\]
			which are natural in $X, Y \in \meas$.
		\item 
			The analogous statement for the distribution functor $\dist : \setc \to \setc$ (resp.~$\distfin : \finset \to \setc$) also holds.
	\end{enumerate}
\end{theo}

For relative product measures, we have to restrict to standard Borel spaces.
But otherwise the statement is the same.

\begin{theo}
	\label{theo:onenaturaltransformation_relative}
	\begin{enumerate}
		\item 
			For any standard Borel space $B$ and $\omega \in \giry B$,
			the formation of relative product measures with respect to $\omega$ is the only family of maps
			\[
				\giry^\omega X \times \giry^\omega Y \longrightarrow \giry^\omega(X \times_B Y)
			\]
			which are natural in $X, Y \in \borelmeas/B$.
		\item 
			With $B$ a (finite) set, the analogous statement for $\dist^\omega : \setc/B \to \setc$ (resp.~$\distfin^\omega : \finset/B \to \setc$) also holds.
	\end{enumerate}
\end{theo}

Taking $B$ to be a singleton recovers the standard Borel version of \cref{theo:onenaturaltransformation}, which naturally holds as well.

\begin{rem}
	\begin{enumerate}
		\item 
			The proofs of both results will first consider the case of finite sets and then extend from there.
			This is designed such that the proofs go through for all of $\giry$, $\dist$ and $\distfin$ simultaneously.
			If one is interested in $\giry$ only,
			then simpler proofs are possible using the Lebesgue measure instead together with the fact that every other probability measure on a standard Borel space is a pushforward of it.
			This is essentially the approach taken in the existing copula theory result~\cite[Theorem~3.1]{Durante2009}.
		\item 
			In \cref{theo:onenaturaltransformation,theo:onenaturaltransformation_relative}, we have considered the functors $\giry$ and $\giry^\omega$ to take values in $\setc$.
			With the corresponding $\giry : \meas \to \meas$ and $\giry^\omega : \borelmeas/B \to \borelmeas$ versions of these functors,
			the same statements hold as obvious corollaries: 
			the only difference is that the natural transformations are now required to have measurable components, which is a stronger condition.
			Therefore our results trivially imply the corresponding statements for these versions.
	\end{enumerate}
\end{rem}

\subsection{Proof of \cref{theo:onenaturaltransformation}}

Although we could prove \cref{theo:onenaturaltransformation} as a special case of \cref{theo:onenaturaltransformation_relative},
we prefer to give a direct proof in order to set the stage for the more general case of relative product measures.
Given a set $X$, we write $\Sym(X)$ for the symmetric group of $X$, which consists of bijections $\sigma: X \rightarrow X$.
We start the proof of \cref{theo:onenaturaltransformation} with the following basic observation.

\begin{lem}\label{lem:uniquepermutationinvariantmeasure}
  Let $X$ be a finite set and $G$ a group acting transitively on $X$. Then, there is exactly one $G$-invariant probability measure on $X$, namely the uniform measure $u_X$.
\end{lem}
\begin{proof}
  This is straightforward based on the fact that a bijection acts on measures by permuting the probabilities of individual outcomes.
\end{proof}

We call $p \in \dist X$ \emph{rational} if $p(x) \in \Q$ for all $x \in X$. A nice property of rational measures is that they are pushforwards of uniform measures on larger sets.
For the precise statement, recall our notation $[N] \coloneqq \{1, \ldots, N\}$.

\begin{lem}\label{lem:rationalmeasure}
  Let $X$ be a finite set and $p \in \dist X$ rational. Then there is $N \in \N$ and $f:\finiteset{N} \rightarrow X$ such that $p = f_\ast u_{\finiteset{N}}$.
\end{lem}
\begin{proof}
  Let $X = \{x_1, \ldots, x_n\}$ and choose a common denominator $N$ for $p(x_1), \ldots, p(x_n)$.
  This means that $p(x_i) = \frac{q_i}{N}$ for some $q_i \in \N$ with $q_1 + \ldots + q_n = N$.
  Define $f: \finiteset{N}  \rightarrow X$ by $f(j) \coloneqq x_{i(j)}$ for every $j \in [N]$,
  where $i(j) \in [n]$ is the unique index with $\sum_{k=1}^{i(j)-1}q_k < j \leq \sum_{k=1}^{i(j)} q_k$.
  Then we have for every $i \in [n]$,
  \[
    f_\ast u_{\finiteset{N}}(x_i) = u_{\finiteset{N}}(f^{-1}(x_i)) =  u_{\finiteset{N}} \left(\left\{\sum_{k=0}^{i-1} q_k + 1,\ldots, \sum_{k=0}^i q_k\right\}\right) = \frac{q_i}{N} = p(x_i).
  \]
  Thus, $p = f_\ast u_{\finiteset{N}}$.
\end{proof}

We can now use these two lemmas to already prove the claim for rational measures.

\begin{prop}\label{prop:rationalcase}
  Let $\dist$ be the distribution functor on $\finset$ and
  \begin{equation}
    \label{eq:alpha}
    \alpha: \dist \_ \times \dist \_ \Rightarrow \dist(\_ \times \_)
  \end{equation}
  any natural transformation.
  Then for all finite sets $X$ and $Y$ and all rational $p \in DX$ and $q \in DY$, we have
  \[
    \alpha(p,q) = p \otimes q.
  \]
\end{prop}

\begin{proof}
  Let $X = \{x_1,\ldots,x_n\}$ and $Y = \{y_1,\ldots,y_m\}$ be finite sets. By naturality of $\alpha$, for all permutations $\sigma \in \operatorname{Sym}(X)$ and $\tau \in \operatorname{Sym}(Y)$,
  \[
    \begin{tikzcd}
      \dist X \times \dist Y \arrow[r, "{\alpha}"] \arrow[d, "\dist\sigma \times \dist\tau"'] & \dist(X\times Y) \arrow[d, "\dist(\sigma \times \tau)"] \\
      \dist X \times \dist Y \arrow[r, "{\alpha}"']                                   & \dist(X \times Y)
    \end{tikzcd}
  \]
  commutes.
  Since the uniform measures $u_X$ and $u_Y$ are invariant under permutations, we have
  \[
    \alpha(u_X,u_Y) = (\sigma \times \tau)_\ast (\alpha(u_X,u_Y)),
  \]
  that is $\alpha(u_X,u_Y)$ is $\operatorname{Sym}(X) \times \operatorname{Sym}(Y)$-invariant. Since the action of $\operatorname{Sym}(X) \times \operatorname{Sym}(Y)$ on $X \times Y$ is transitive, we conclude its uniformity
  \[
	  \alpha(u_X,u_Y) = u_{X \times Y}
  \]
  by \cref{lem:uniquepermutationinvariantmeasure}.

  Now let $p \in \dist X$ and $q \in \dist Y$ be rational.
  Then \cref{lem:rationalmeasure} provides us with $N,M \in \N$ and maps $f:\finiteset{N} \rightarrow X$ and $g:\finiteset{M} \rightarrow Y$ such that $p = f_\ast u_{\finiteset{N}}$ and $q = g_\ast u_{\finiteset{M}}$. 
  %In particular, $p(x_i) = \frac{A_i}{N}$ and $q(y_j) = \frac{B_j}{M}$ for all $i,j$ and $A_i,B_j,M,N \in \N$.
  Thus, commutativity of
  \[
    \begin{tikzcd}
      \dist \finiteset{N} \times \dist \finiteset{M} \arrow[r, "{\alpha}"] \arrow[d, "\dist f \times \dist g"'] & \dist(\finiteset{N} \times \finiteset{M}) \arrow[d, "\dist(f \times g)"] \\
      \dist X \times \dist Y \arrow[r, "{\alpha}"']                               & \dist(X \times Y)
    \end{tikzcd}
  \]
  gives the second step in
  \begin{align*}
	  \alpha(p,q) & = \alpha(f_\ast u_{\finiteset{N}},g_\ast u_{\finiteset{M}}) \\[2pt]
			    & = (f \times g)_\ast u_{\finiteset{N} \times \finiteset{M}} \\[2pt]
			    & = (f \times g)_\ast (u_{\finiteset{N}} \otimes u_{\finiteset{M}}) \\[2pt]
			    & = f_\ast u_{\finiteset{N}} \otimes g_\ast u_{\finiteset{M}} \\[2pt]
			    & = p \otimes q,
  \end{align*}
  while the fourth step is the naturality of the product measure construction $\nabla$.
\end{proof}

The general case without the restriction to rational measures can now be obtained by a density argument. The elegance of this argument lies in translating the natural order on the interval $[0,1]$ into a purely probabilistic condition. This is very much in the spirit of \emph{effectus theory}, a categorical framework for probabilistic logic \cite{Jacobs2015}.

\begin{lem}\label{prop:order}
  Define $\theta^-,\theta^+: \three \rightarrow \two$ by
  \begin{align*}
	  \theta^-(1) = \theta^-(2) = 1, &&  \theta^+(1) = 1,  \\
	  \theta^-(3) = 2, && \theta^+(2) = \theta^+(3) = 2.
  \end{align*}
  Then for all $p, q \in \dist\two$, we have $p(1) \leq q(1)$ if and only if there is $\omega \in \dist\three$ such that
  \begin{equation}\label{eq:theta}
    \theta^+_\ast\omega = p \qquad \text{and} \qquad \theta^-_\ast\omega = q.
  \end{equation}
  In this case, the only such $\omega$ satisfies
  \begin{equation}
	  \label{eq:theta2}
	  \omega(1) = p(1), \qquad \omega(2) = q(1)-p(1), \qquad \omega(3) = q(2).
  \end{equation}
\end{lem}
\begin{proof}
  It is straightforward that the only solution of the linear system \cref{eq:theta} is given by \cref{eq:theta2} and that it satisfies the normalization condition $\omega(1) + \omega(2) + \omega(3) = 1$.
  It is nonnegative if and only if $p(1) \leq q(1)$.
\end{proof}

If we identify $D\two \cong [0,1]$ via $p \mapsto p(1)$, the relation $p(1) \leq q(1)$ becomes the natural order on $[0,1]$.
Hence  the natural order on $[0,1]$ induces an order on $\dist\two$ which is witnessed, in purely probabilistic terms, by $\omega$.

\begin{lem}\label{lem:bernoulliirrational}
	For all $p, q \in \dist\two$, every natural transformation $\alpha$ as in~\eqref{eq:alpha} satisfies
  \[
	  \alpha(p,q) = p \otimes q.
  \]
\end{lem}
\begin{proof}
	We first treat the case where $q$ is rational.
  For $T \subseteq \two$, define
  \begin{align*}
		\phi_T \: : \: \dist\two & \longrightarrow [0,1] \\
		s & \longmapsto \alpha(s,q)(\{1\} \times T).
  \end{align*}
	Let $p,p' \in \dist \two$ satisfy $p \leq p'$, and let $\omega \in \dist\three$ be the witness of this inequality as per \cref{prop:order}.
	Then,
  \begin{align*}
    \phi_T(p') - \phi_T(p) &= ((\theta^- \times \id)_\ast\alpha(\omega,q))(\{1\} \times T) - ((\theta^+ \times \id)_\ast\alpha(\omega,q))(\{1\} \times T) = \\
                          &= \alpha(\omega,q)(\{1,2\} \times T) - \alpha(\omega,q)(\{1\} \times T) = \\
                          &= \alpha(\omega,q)(\{2\} \times T) \geq 0.
  \end{align*}
  Hence $\phi_T$ is monotonically nondecreasing.

	Let $p \in \dist\two$ be arbitrary. By density of $\Q$ in $\R$, we can find rational $(p_n^-)_{n \in \N}$ and $(p_n^+)_{n \in \N}$ in $\dist\two$ 
	with $p_n^- \nearrow p$ and $p^+_n \searrow p$. Thus, by monotonicity,
    \[
      \phi_T(p^-_n) \leq \phi_T(p) \leq \phi_T(p^+_n)
    \]
  for all $n \in \N$. By definition of $\phi_T$ and \cref{prop:rationalcase}, we thus obtain
  \[
    p^-_n(1) \cdot q(T) \leq \alpha(p,q)(\{1\} \times T) \leq p^+_n(1) \cdot q(T).
  \]
  Taking $n \to \infty$ shows that $\alpha(p,q)(\{1\} \times T) = p(1) \cdot q(T)$ for all $T \subseteq \two$.
	Since we can make the same argument with $\{2\}$ in place of $\{1\}$, the claim follows by the uniqueness part of \cref{defi:productmeasure}.

	The statement is now proven for all $p$ and rational $q$.
	To lift to arbitrary $q$, we can use the same rational approximation argument as above, but now with respect to $q$ instead of $p$.
\end{proof}

The final steps are now relatively straightforward.

\begin{proof}[Proof of \cref{theo:onenaturaltransformation}]
	Let $\alpha : \giry\_ \times \giry\_ \Rightarrow \giry(\_ \times \_)$ be any natural transformation.
	By \cref{lem:bernoulliirrational}, $\alpha$ and $\nabla$ agree on $\giry\two \times \giry\two$.
  If $X$ and $Y$ are arbitrary measurable spaces, let $p \in \giry(X)$ and $q \in \giry(Y)$ be given.
  For any $S \in \Sigma_X$ and $T \in \Sigma_Y$, consider the pushforwards of $p$ and $q$ along the indicator functions,
	\[
		p_S \coloneq \chi^S_\ast p, \qquad q_T \coloneq \chi^T_\ast q.
	\]
	Then we already know $\alpha(p_S,q_T) = p_S \otimes q_T$ by the finite case.
	Therefore, and by naturality of $\alpha$,
	\begin{align*}
		\alpha(p,q)(S \times T) & = ((\chi^S \times \chi^T)_\ast\alpha(p,q))(\{1\} \times \{1\}) \\[2pt]
																		& = \alpha(p_S,q_T)(\{1\} \times \{1\}) \\[2pt]
																	& = p(S) \cdot q(T).
	\end{align*}
	The claimed $\alpha = \nabla$ hence follows as before.

	By restricting the proof from $\giry$ to $\dist$ and $\distfin$, we also obtain the analogous statement for the distribution functor on $\setc$.
\end{proof}

\subsection{Proof of \cref{theo:onenaturaltransformation_relative}}

We follow the same steps as for ordinary product measures: first measures with uniform fibers, then rational probabilities,
then an order-theoretic approximation argument on $\two$, and finally indicator functions.

Throughout, fix the base space $B \in \borelmeas$, the base measure $\omega \in \giry B$ and a natural transformation
\[
	\alpha \: : \: \giry^\omega\_ \times \giry^\omega\_ \Longrightarrow \giry^\omega(\_ \times_B \_).
\]
Our goal is to prove that $\alpha = \nabla^\omega$.

For a finite set $X$ and a Markov kernel $p : B \to \dist X$, write
\[
	\overline{p} \coloneq \int_{b \in B} p(-|b) \otimes \delta_b \, \omega(\dx b)
	\quad\in\quad \giry^\omega(X \times B),
\]
where $X \times B$ is regarded as an object over $B$ via the projection.
For example with $u_X$ the uniform measure on $X$, and with every measure regarded as a constant kernel, 
we get $\overline{u_X} = u_X \otimes \omega$.
In general, we have
\[
	\overline{p}(\{x\} \times S) = \int_{b \in S} p(x|b)\,\omega(\dx b) \qquad \forall S \in \Sigma_B.
\]
Hence the kernel $b \mapsto p(-|b) \otimes \delta_b$ is a regular conditional probability of $\overline{p}$ with respect to the projection to $B$.

The analogue of rational measures that we need consists of kernels whose probabilities have a common denominator over the whole base:

\begin{prop}\label{prop:relative_rationalcase}
	Let $X,Y$ be nonempty finite sets and $p : B \to \dist X$, $q : B \to \dist Y$ Markov kernels.
	Suppose that there are positive integers $N,M$ such that
	\[
		Np(x|b) \in \N, \qquad Mq(y|b) \in \N
		\qquad\quad
		\forall x \in X, y \in Y, b \in B.
	\]
	Then
	\[
		\alpha(\overline{p},\overline{q}) = \overline{p} \otimes_\omega \overline{q}.
	\]
\end{prop}

From now on, we will freely use the identification
\[
	(X \times B) \times_B (Y \times B) \cong X \times Y \times B,
\]
which frequently allows for more convenient notation in the proofs.

\begin{proof}
	We first consider the uniform measures $u_X$ and $u_Y$.
	Naturality under $\sigma \times \id_B$ and $\tau \times \id_B$, for permutations $\sigma$ of $X$ and $\tau$ of $Y$, implies that
	\[
		r \coloneq \alpha(\overline{u_X},\overline{u_Y})
	\]
	is invariant under $(\sigma \times \tau) \times \id_B$.
	For every $S \in \Sigma_B$, transitivity of the action on $X \times Y$ therefore makes $r(\{(x,y)\} \times S)$ independent of $(x,y)$.
	Since $r$ has pushforward $\omega$ on $B$,
	the only possible value is
	\[
		r(\{(x,y)\} \times S) = \frac{\omega(S)}{|X|\,|Y|}.
	\]
	Thus $r = u_{X \times Y} \otimes \omega = \overline{u_X} \otimes_\omega \overline{u_Y}$.

	Now let $p,q$ satisfy the stated denominator conditions.
	As in \cref{lem:rationalmeasure}, write $X = \{x_1,\ldots,x_n\}$ and define a map over $B$ by
	\begin{align*}
		f \: : \: \finiteset{N} \times B & \longrightarrow X \times B \\
		(j,b) & \longmapsto (x_i,b),
	\end{align*}
	where $x_i$ is the unique element of $X$ with
	\[
		\sum_{k=1}^{i-1} Np(x_k|b) < j \leq \sum_{k=1}^{i} Np(x_k|b).
	\]
	The defining inequalities are measurable in $b$, and therefore $f$ is measurable.
	For each $b$, exactly $Np(x_i|b)$ elements map to $(x_i,b)$, and hence $f_\ast\overline{u_{\finiteset{N}}} = \overline{p}$.
	Similarly, there is a measurable map $g : \finiteset{M} \times B \to Y \times B$ over $B$ with $g_\ast\overline{u_{\finiteset{M}}} = \overline{q}$.
	By naturality, the uniform case, and \cref{lem:relativeproductnaturality},
	we have
	\begin{align*}
		\alpha(\overline{p},\overline{q})
		&= (f \times_B g)_\ast
		\alpha
		(\overline{u_{\finiteset{N}}},\overline{u_{\finiteset{M}}}) \\[2pt]
		&= (f \times_B g)_\ast
		\bigl(\overline{u_{\finiteset{N}}} \otimes_\omega \overline{u_{\finiteset{M}}}\bigr) \\[2pt]
		&= \overline{p} \otimes_\omega \overline{q}.
		\qedhere
	\end{align*}
\end{proof}

We next extend from these rational kernels to arbitrary Bernoulli kernels, in parallel with \cref{lem:bernoulliirrational}.
This means that for kernels $p,p' : B \to \dist\two$, we write
\[
	p \le p' \qquad \Longleftrightarrow \qquad
	p(1|b) \leq p'(1|b) \quad \forall b \in B.
\]
The witness in \cref{prop:order} can be chosen measurably: define $\eta : B \to \dist\three$ by
\[
	\eta(1|b) = p(1|b), \qquad
	\eta(2|b) = p'(1|b)-p(1|b), \qquad
	\eta(3|b) = p'(2|b).
\]
Writing $\theta_B^\pm \coloneq \theta^\pm \times \id_B$ then gives
\[
	(\theta_B^+)_\ast\overline{\eta} = \overline{p},
	\qquad
	(\theta_B^-)_\ast\overline{\eta} = \overline{p}'.
\]

\begin{lem}\label{lem:relative_bernoulliirrational}
	For all Markov kernels $p,q : B \to \dist\two$, we have
	\[
		\alpha(\overline{p},\overline{q})
		= \overline{p} \otimes_\omega \overline{q}.
	\]
\end{lem}

\begin{proof}
	We first assume that $q$ has a common denominator as in \cref{prop:relative_rationalcase}.
	For $T \subseteq \two$ and $S \in \Sigma_B$, define
	\[
		\phi_{T,S}(s) \coloneq
		\alpha(\overline{s},\overline{q})(\{1\} \times T \times S).
	\]
	If $p \leq p'$, let $\eta$ be the witness above and put
	$r \coloneq \alpha(\overline{\eta},\overline{q})$.
	Naturality under $\theta_B^\pm$ gives
	\begin{align*}
		\phi_{T,S}(p')-\phi_{T,S}(p)
		&= r(\{1,2\} \times T \times S)-r(\{1\} \times T \times S) \\
		&= r(\{2\} \times T \times S) \geq 0.
	\end{align*}
	Thus $\phi_{T,S}$ is monotonically nondecreasing.

	For an arbitrary kernel $p$, define measurable dyadic approximations by
	\[
		p_n^-(1|b) \coloneq 2^{-n}\lfloor 2^n p(1|b)\rfloor,
		\qquad
		p_n^+(1|b) \coloneq 2^{-n}\lceil 2^n p(1|b)\rceil,
	\]
	and $p_n^\pm(2|b) \coloneq 1-p_n^\pm(1|b)$.
	These satisfy $p_n^- \nearrow p$ and $p_n^+ \searrow p$ pointwise, and both have common denominator $2^n$.
	By monotonicity and \cref{prop:relative_rationalcase},
	\[
		\int_S p_n^-(1|b) \, q(T|b) \, \omega(\dx b)
		\leq \phi_{T,S}(p)
		\leq \int_S p_n^+(1|b) \, q(T|b)\,\omega(\dx b).
	\]
	Taking $n \to \infty$ and using dominated convergence gives
	\[
		\phi_{T,S}(p) = \int_S p(1|b) \, q(T|b) \, \omega(\dx b).
	\]
	We likewise get the same identity with $\{2\}$ in place of $\{1\}$.
	Since evaluation on sets of the form $\{i\} \times \{j\} \times S$ distinguishes measures on $\two \times \two \times B$,
	the claim follows for arbitrary $p$ whenever $q$ has a common denominator.

	Now fix arbitrary $p$ and $q$.
	Applying the same witness construction in the second variable shows that, for every $T \subseteq \two$ and $S \in \Sigma_B$, the map
	\begin{align*}
		\dist\two & \longrightarrow [0,1] \\
		t & \longmapsto \alpha(\overline{p},\overline{t})(T \times \{1\} \times S)
	\end{align*}
	is monotonically nondecreasing.
	Approximate $q$ from below and above by $q_n^-$ and $q_n^+$ as above.
	The case already proved and dominated convergence give
	\[
		\alpha(\overline{p},\overline{q})(T \times \{1\} \times S)
		= \int_S p(T|b) \, q(1|b) \, \omega(\dx b).
	\]
	Again we get the same identity with $\{2\}$ in place of $\{1\}$ as well, proving the claim.
\end{proof}

The final reduction again uses indicator functions, now paired with the structure maps to $B$.

\begin{proof}[Proof of \cref{theo:onenaturaltransformation_relative}]
	Given arbitrary $X,Y \in \borelmeas/B$,
	we show that $\alpha(p,q) = p \otimes_\omega q$ for all $p \in \giry^\omega X$ and $q \in \giry^\omega Y$.
	Choose regular conditional probabilities $p(-|b)$ and $q(-|b)$ using \cref{prop:regularconditionalprobability}.
	For $S \in \Sigma_X$ and $T \in \Sigma_Y$, define measurable maps $f_S : X \to \two \times B$ and $g_T : Y \to \two \times B$ over $B$ by
	\[
		f_S(x) \coloneq (\chi_S(x),\rho_X(x)),
		\qquad
		g_T(y) \coloneq (\chi_T(y),\rho_Y(y)),
	\]
	where $\chi_S$ and $\chi_T$ are the respective indicator functions.
	Also consider the Markov kernels $p_S,q_T : B \to \dist\two$ characterized by
	\[
		p_S(\{1\}|b) \coloneq p(S|b), \qquad q_T(\{1\}|b) \coloneq q(T|b).
	\]
	The disintegration identity~\eqref{eq:disintegration} gives
	$(f_S)_\ast p = \overline{p_S}$ and $(g_T)_\ast q = \overline{q_T}$.
	Naturality and \cref{lem:relative_bernoulliirrational} therefore imply
	\begin{align*}
		\alpha(p,q)\bigl((S \times T) \cap (X \times_B Y)\bigr)
		&= \bigl((f_S \times_B g_T)_\ast\alpha(p,q)\bigr)(\{1\} \times \{1\} \times B) \\[2pt]
		&= \alpha(\overline{p_S},\overline{q_T})(\{1\} \times \{1\} \times B) \\[2pt]
		&= \int_B p(S|b) \, q(T|b) \, \omega(\dx b) \\[2pt]
		&= (p \otimes_\omega q)\bigl((S \times T) \cap (X \times_B Y)\bigr).
	\end{align*}
	The restrictions of measurable rectangles form a generating $\pi$-system for $X \times_B Y$.
	Thus $\alpha(p,q) = p \otimes_\omega q$, as required.

	Again, the restriction of the whole proof to $\dist^\omega$ and $\distfin^\omega$ is straightforward.
\end{proof}





\printbibliography[heading=bibintoc]
\end{document}

\typeout{get arXiv to do 4 passes: Label(s) may have changed. Rerun}
%%% Local Variables:
%%% mode: LaTeX
%%% TeX-master: t
%%% End':
